\chapter{2004年重庆卷（文）}
\section{单选题}
\begin{problem}
函数 \( y = \sqrt{\log_{\frac{1}{2}}(3x - 2)} \) 的定义域是（\quad）

\choices
    {\([1, +\infty)\)}
    {\(\left(\frac{2}{3}, +\infty\right)\)}
    {\(\left[\frac{2}{3}, 1\right]\)}
    {\(\left(\frac{2}{3}, 1\right]\)}
\end{problem}

\begin{answer}
D
\end{answer}

\begin{solution}
要使根式有意义，必须满足 \(\log_{\frac{1}{2}}(3x-2)\geqslant0\). 因为对数的真数还需为正，且底数 \(\frac{1}{2}\) 小于 \(1\)，所以 \(0<3x-2\leqslant1\).
解得 \(\frac{2}{3}<x\leqslant1\)，故定义域为 \(\left(\frac{2}{3},1\right]\)，选 D.
\end{solution}

\begin{problem}
函数 \( f(x) = \frac{x^2 - 1}{x^2 + 1} \)，则 \( \frac{f(2)}{f\left(\frac{1}{2}\right)} = \)（\quad）

\choices
    {\(1\)}
    {\(-1\)}
    {\(\frac{3}{5}\)}
    {\(-\frac{3}{5}\)}
\end{problem}

\begin{answer}
B
\end{answer}

\begin{solution}
代入得 \(f(2)=\frac{2^2-1}{2^2+1}=\frac{3}{5}\)，\(f\left(\frac{1}{2}\right)=\frac{\frac{1}{4}-1}{\frac{1}{4}+1}=-\frac{3}{5}\). 因此 \(\frac{f(2)}{f\left(\frac{1}{2}\right)}=\frac{\frac{3}{5}}{-\frac{3}{5}}=-1\)，
故选 B.
\end{solution}

\begin{problem}
圆 \(x^{2} + y^{2} - 2x + 4y + 3 = 0\) 的圆心到直线 \(x - y = 1\) 的距离为（\quad）

\choices
    {\(2\)}
    {\(\frac{\sqrt{2}}{2}\)}
    {\(1\)}
    {\(\sqrt{2}\)}
\end{problem}

\begin{answer}
D
\end{answer}

\begin{solution}
将圆的方程配方，得 \((x-1)^2+(y+2)^2=2\)，所以圆心为 \(C(1,-2)\). 圆心到直线 \(x-y=1\)，即 \(x-y-1=0\) 的距离为 \(\frac{|1-(-2)-1|}{\sqrt{1^2+(-1)^2}}=\frac{2}{\sqrt{2}}=\sqrt{2}\)，
故选 D.
\end{solution}

\begin{problem}
不等式 \(x + \frac{2}{x + 1} > 2\) 的解集是（\quad）

\choices
    {\((-1,0)\cup (1, + \infty)\)}
    {\((-\infty, -1) \cup (0, 1)\)}
    {\((-1,0)\cup (0,1)\)}
    {\((-\infty, -1) \cup (1, + \infty)\)}
\end{problem}

\begin{answer}
A
\end{answer}

\begin{solution}
原不等式的定义域为 \(x\neq-1\). 移项通分，得 \(x+\frac{2}{x+1}-2=\frac{x(x-1)}{x+1}>0\). 分式的分界点为 \(-1\)、\(0\)、\(1\)，将数轴分成四个区间：在 \((- \infty,-1)\) 上分子为正、分母为负，分式为负；在 \((-1,0)\) 上分子、分母均为正，分式为正；在 \((0,1)\) 上分子为负、分母为正，分式为负；在 \((1,+\infty)\) 上分子、分母均为正，分式为正. 又因不等式为严格不等式，\(x=0\)、\(x=1\) 不取，且 \(-1\) 不在定义域内，所以解集为 \((-1,0)\cup(1,+\infty)\)，故选 A.
\end{solution}

\begin{problem}
\(\sin 163^{\circ}\sin 223^{\circ} + \sin 253^{\circ}\sin 313^{\circ} =\)（\quad）

\choices
    {\(-\frac{1}{2}\)}
    {\(\frac{1}{2}\)}
    {\(-\frac{\sqrt{3}}{2}\)}
    {\(\frac{\sqrt{3}}{2}\)}
\end{problem}

\begin{answer}
B
\end{answer}

\begin{solution}
利用积化和差公式，原式为 \(\frac{1}{2}\left(\cos60^\circ-\cos386^\circ\right)+\frac{1}{2}\left(\cos60^\circ-\cos566^\circ\right)\). 又 \(\cos386^\circ=\cos26^\circ\)，\(\cos566^\circ=\cos206^\circ=-\cos26^\circ\)，所以原式为 \(\frac{1}{2}\left(\frac{1}{2}-\cos26^\circ\right)+\frac{1}{2}\left(\frac{1}{2}+\cos26^\circ\right)=\frac{1}{2}\)，
故选 B.
\end{solution}

\begin{problem}
若向量 \(\bs{a}\) 与 \(\bs{b}\) 的夹角为 \(60^{\circ}\)，\(\left|\bs{b}\right| = 4\)，\(\left(\bs{a} + 2\bs{b}\right) \cdot \left(\bs{a} - 3\bs{b}\right) = -72\)，则向量 \(\bs{a}\) 的模为（\quad）

\choices
    {\(2\)}
    {\(4\)}
    {\(6\)}
    {\(12\)}
\end{problem}

\begin{answer}
C
\end{answer}

\begin{solution}
设 \(\left|\bs{a}\right|=t\)，则 \(\bs{a}\cdot\bs{b}=\left|\bs{a}\right|\left|\bs{b}\right|\cos60^\circ=2t\). 将已知等式左边展开，得 \((\bs{a}+2\bs{b})\cdot(\bs{a}-3\bs{b})=|\bs{a}|^2-\bs{a}\cdot\bs{b}-6|\bs{b}|^2=t^2-2t-96\). 因此 \(t^2-2t-96=-72\)，即 \(t^2-2t-24=0\)，
解得 \(t=6\) 或 \(t=-4\). 因为模非负，所以 \(t=6\)，故选 C.
\end{solution}

\begin{problem}
已知 \( p \) 是 \( r \) 的充分不必要条件，\( s \) 是 \( r \) 的必要条件，\( q \) 是 \( s \) 的必要条件，那么 \( p \) 是 \( q \) 成立的（\quad）

\choices
    {充分不必要条件}
    {必要不充分条件}
    {充要条件}
    {既不充分也不必要条件}
\end{problem}

\begin{answer}
A
\end{answer}

\begin{solution}
由“\(p\) 是 \(r\) 的充分不必要条件”知 \(p\Rightarrow r\)，且 \(r\) 不推出 \(p\). 由“\(s\) 是 \(r\) 的必要条件”知 \(r\Rightarrow s\)，由“\(q\) 是 \(s\) 的必要条件”知 \(s\Rightarrow q\). 因而
\(p\Rightarrow r\Rightarrow s\Rightarrow q\).
所以 \(p\) 是 \(q\) 的充分条件. 又存在 \(r\) 成立而 \(p\) 不成立的情形，此时由 \(r\Rightarrow s\Rightarrow q\) 可知 \(q\) 成立而 \(p\) 不成立，因此 \(p\) 不是 \(q\) 的必要条件. 故 \(p\) 是 \(q\) 成立的充分不必要条件，选 A.
\end{solution}

\begin{problem}
不同直线 \(m\)，\(n\) 和不同平面 \(\alpha\)，\(\beta\)，给出下列命题：

\begin{circlelist}
\item \(\begin{cases}
\alpha&\parallel\beta,\\
m&\subset\alpha,
\end{cases}\Rightarrow m\parallel\beta\)；
\item \(\begin{cases}
m&\parallel n,\\
m&\parallel\beta,
\end{cases}\Rightarrow n\parallel\beta\)；
\item \(\begin{cases}
m&\subset\alpha,\\
n&\subset\beta,
\end{cases}\Rightarrow m\)，\(n\) 异面；
\item \(\begin{cases}
\alpha&\perp\beta,\\
m&\parallel\alpha,
\end{cases}\Rightarrow m\perp\beta\).
\end{circlelist}

其中假命题有（\quad）

\choices
    {\(0\) 个}
    {\(1\) 个}
    {\(2\) 个}
    {\(3\) 个}
\end{problem}

\begin{answer}
D
\end{answer}

\begin{solution}
第一个命题正确. 因为 \(\alpha\parallel\beta\)，平面 \(\alpha\) 内的任意直线都不与平面 \(\beta\) 相交；又 \(m\subset\alpha\)，所以 \(m\parallel\beta\).

第二个命题不一定正确. 在平面 \(\beta\) 内取一条直线 \(n\)，在与 \(\beta\) 平行的另一平面内取一条与 \(n\) 平行的直线 \(m\). 则 \(m\parallel n\) 且 \(m\parallel\beta\)，但 \(n\subset\beta\)，按直线与平面平行的定义，\(n\) 不平行于 \(\beta\). 因而第二个命题是假命题.

第三个命题不一定正确. 若 \(\alpha\) 与 \(\beta\) 相交，设交线为 \(l\)，在 \(\alpha\) 和 \(\beta\) 内分别取经过 \(l\) 上同一点的两条不同直线 \(m\)、\(n\)，则 \(m\) 与 \(n\) 相交，不是异面直线. 因而第三个命题是假命题.

第四个命题不一定正确. 设 \(l=\alpha\cap\beta\)，取一条与 \(l\) 平行且不在 \(\alpha\) 内的直线 \(m\). 由于 \(l\subset\alpha\)，有 \(m\parallel\alpha\)；但 \(m\) 的方向与平面 \(\beta\) 内的直线 \(l\) 平行，且可取为不与 \(\beta\) 相交，所以不可能垂直于平面 \(\beta\). 因而第四个命题是假命题.

所以假命题有 3 个，故选 D.
\end{solution}

\begin{problem}
若数列 \( \{a_{n}\} \) 是等差数列，首项 \(a_{1}>0\)，\(a_{2003}+a_{2004}>0\)，\(a_{2003}\cdot a_{2004}<0\)，则使前 \(n\) 项和 \( S_{n}>0 \) 成立的最大自然数 \(n\) 是（\quad）

\choices
    {\(4005\)}
    {\(4006\)}
    {\(4007\)}
    {\(4008\)}
\end{problem}

\begin{answer}
B
\end{answer}

\begin{solution}
若 \(a_{2003}<0<a_{2004}\)，则公差 \(d=a_{2004}-a_{2003}>0\)，从而 \(a_1=a_{2003}-2002d<0\)，与 \(a_1>0\) 矛盾. 因此 \(a_{2003}>0>a_{2004}\). 结合 \(a_{2003}a_{2004}<0\) 和 \(a_{2003}+a_{2004}>0\)，设
\(a_{2003}=u>0\)，\(a_{2004}=-v<0\)
其中 \(u>v>0\). 设等差数列的公差为 \(d\)，则
\[
d=a_{2004}-a_{2003}=-(u+v)
\]
并且
\[
a_1=a_{2003}-2002d=2003u+2002v
\]
等差数列前 \(n\) 项和为
\[
S_n=\frac{n}{2}(a_1+a_n)
\]
因为 \(n>0\)，所以
\[
S_n>0\Longleftrightarrow a_1+a_n>0\Longleftrightarrow 2a_1+(n-1)d>0
\]
代入 \(a_1\) 和 \(d\)，得
\[
n-1<\frac{2(2003u+2002v)}{u+v}=4004+\frac{2u}{u+v}
\]
由 \(u>v>0\) 知 \(\frac{1}{2}<\frac{u}{u+v}<1\)，从而
\[
4005<\frac{2(2003u+2002v)}{u+v}<4006
\]
因此满足不等式的最大整数 \(n-1\) 为 \(4005\)，故 \(n\) 的最大值为 \(4006\)，选 B.
\end{solution}

\begin{problem}
已知双曲线 \(\frac{x^2}{a^2} - \frac{y^2}{b^2} = 1\)（\(a > 0\)，\(b > 0\)）的左，右焦点分别为 \(F_1\)，\(F_2\)，点 \(P\) 在双曲线的右支上，且 \(|PF_1| = 4|PF_2|\). 则双曲线的离心率 \(e\) 的最大值为（\quad）

\choices
    {\( \frac{4}{3} \)}
    {\( \frac{5}{3} \)}
    {\(2\)}
    {\( \frac{7}{3} \)}
\end{problem}

\begin{answer}
B
\end{answer}

\begin{solution}
设双曲线的半焦距为 \(c\)，则离心率 \(e=\frac{c}{a}\). 对双曲线右支上的点 \(P\)，有双曲线的焦点性质
\[
|PF_1|-|PF_2|=2a
\]
结合 \(|PF_1|=4|PF_2|\)，得
\(3|PF_2|=2a\)，\(|PF_2|=\frac{2a}{3}\).
右支上的点可写为 \(P=(x,y)\)，其中 \(x\geqslant a\)，且
\[
y^2=\frac{b^2}{a^2}(x^2-a^2)
\]
因此
\[
|PF_2|^2=(x-c)^2+y^2=\frac{c^2}{a^2}x^2-2cx+a^2=\left(\frac{cx}{a}-a\right)^2
\]
由于 \(x\geqslant a\) 且 \(c>a\)，有 \(|PF_2|=\frac{cx}{a}-a\geqslant c-a\)，等号当且仅当 \(P\) 为右顶点. 所以
\[
c-a\leqslant\frac{2a}{3}
\]
即 \(c\leqslant\frac{5a}{3}\). 因而
\[
e=\frac{c}{a}\leqslant\frac{5}{3}
\]
当 \(P\) 取右顶点且 \(c=\frac{5a}{3}\) 时，\(|PF_2|=c-a=\frac{2a}{3}\)，\(|PF_1|=c+a=\frac{8a}{3}=4|PF_2|\)，题设条件确实成立，所以等号可以取得，离心率的最大值为 \(\frac{5}{3}\)，选 B.
\end{solution}

\begin{problem}
已知盒中装有 \(3\) 只螺口与 \(7\) 只卡口灯泡，这些灯泡的外形与功率都相同且灯口向下放着，现需要一只卡口灯泡使用，电工师傅每次从中任取一只并不放回，则他直到第 \(3\) 次才取得卡口灯泡的概率为（\quad）

\choices
    {\( \frac{21}{40} \)}
    {\( \frac{17}{40} \)}
    {\( \frac{3}{10} \)}
    {\( \frac{7}{120} \)}
\end{problem}

\begin{answer}
D
\end{answer}

\begin{solution}
要使第 3 次取到卡口灯泡，前两次必须取到螺口灯泡，且第 3 次取到卡口灯泡. 因此所求概率为
\[
\frac{3}{10}\cdot\frac{2}{9}\cdot\frac{7}{8}=\frac{7}{120}
\]
故选 D.
\end{solution}

\begin{problem}
如图，棱长为 \(5\) 的正方体无论从哪一个面看，都有两个直通的边长为 \(1\) 的正方形孔. 则这个有孔正方体的表面积（含孔内各面）是（\quad）

\bitmapfigure[max width=0.46\linewidth]{img/2004/chongqing_liberal/q12_fig1.png}

\choices
    {\(258\)}
    {\(234\)}
    {\(222\)}
    {\(210\)}
\end{problem}

\begin{answer}
C
\end{answer}

\begin{solution}
正方体六个外表面上各有两个边长为 \(1\) 的正方形孔口，所以外表面实际保留的面积为
\[
S_{\text{外}}=6\left(5^2-2\cdot1^2\right)=138
\]
六个直通孔可看成六条长为 \(5\)、截面为边长 \(1\) 的正方形通道，每条通道沿长度方向由五个单位小立方体组成. 由图可知，每条通道有两个位置与其他通道相交；六条通道的交会次数按两个通道共用一个交会小立方体计算，因此共有 \(6\times2\div2=6\) 个交会小立方体. 除去这些交会位置后，每条通道还有三个不交会的小立方体，每个贡献四个单位内壁面，所以这部分面积为 \(6\times3\times4=72\). 每个交会小立方体由两条互相垂直的孔道共用，按图中实际边界计数只留下两个单位内壁面，故交会处的面积为 \(6\times2=12\). 因而孔内各面的总面积为 \(S_{\text{内}}=72+12=84\).

所以有孔正方体的表面积为 \(S=S_{\text{外}}+S_{\text{内}}=138+84=222\)，故选 C.
\end{solution}

\section{填空题}
\begin{problem}
若在 \((1 + a x)^5\) 的展开式中 \(x^3\) 的系数为 \(-80\)，则 \(a=\)\fillinblank{}.
\end{problem}

\begin{answer}
\(-2\)
\end{answer}

\begin{solution}
由二项式定理，\((1+ax)^5\) 中 \(x^3\) 的系数为 \(\mathrm C_5^3a^3=10a^3\). 因此 \(10a^3=-80\)，即 \(a^3=-8\)，所以 \(a=-2\).
\end{solution}

\begin{problem}
已知 \( \frac{5}{x} + \frac{3}{y} = 2 \)（\(x > 0\)，\(y > 0\)），则 \(xy\) 的最小值是\fillinblank{}.
\end{problem}

\begin{answer}
\(15\)
\end{answer}

\begin{solution}
令 \(u=\frac{5}{x}\)，\(v=\frac{3}{y}\)，则 \(u>0\)，\(v>0\)，且 \(u+v=2\). 由 \((u-v)^2\geqslant0\)，得 \(uv\leqslant1\). 又 \(xy=\frac{15}{uv}\)，所以 \(xy\geqslant15\). 当且仅当 \(u=v=1\)，即 \(x=5\)，\(y=3\) 时取等号，故 \(xy\) 的最小值为 \(15\).
\end{solution}

\begin{problem}
已知曲线 \( y = \frac{1}{3} x^3 + \frac{4}{3} \)，则过点 \( P(2,4) \) 的切线方程是\fillinblank{}.
\end{problem}

\begin{answer}
\(y-4x+4=0\)
\end{answer}

\begin{solution}
对曲线方程求导，得 \(y'=x^2\). 在点 \(P(2,4)\) 处，切线斜率为 \(2^2=4\)，因而切线方程为 \(y-4=4(x-2)\)，整理得 \(y-4x+4=0\).
\end{solution}

\begin{problem}
毛泽东在《送瘟神》中写到：“坐地日行八万里”. 又知地球的体积大约是火星的 \(8\) 倍，则火星的大圆周长约\fillinblank{}\(\text{万里}\).
\end{problem}

\begin{answer}
\(4\)
\end{answer}

\begin{solution}
设地球和火星的半径分别为 \(R\) 和 \(r\). 由球的体积公式及题意，得 \(\frac{4}{3}\pi R^3=8\cdot\frac{4}{3}\pi r^3\)，所以 \(R^3=8r^3\)，从而 \(R=2r\). 因此地球与火星的大圆周长之比为 \(\frac{2\pi R}{2\pi r}=2\). 地球大圆周长约为八万里，所以火星大圆周长约为 \(8\div2=4\)\(\text{万里}\).
\end{solution}

\section{解答题}
\begin{problem}
求函数 \( y = \sin^4 x + 2\sqrt{3}\sin x\cos x - \cos^4 x \) 的最小正周期和最小值；并写出该函数在 \([0, \pi]\) 上的单调递增区间.
\end{problem}

\begin{solution}
利用 \(\sin^4x-\cos^4x=(\sin^2x-\cos^2x)(\sin^2x+\cos^2x)=-\cos2x\)，以及 \(2\sin x\cos x=\sin2x\)，得
\[
y=\sqrt{3}\sin2x-\cos2x=2\sin\left(2x-\frac{\pi}{6}\right)
\]
由于正弦函数的最小正周期为 \(2\pi\)，而自变量的系数为 \(2\)，所以函数的最小正周期为 \(\pi\)，最小值为 \(-2\).

函数的导数为 \(y'=2\sqrt{3}\cos2x+2\sin2x=4\cos\left(2x-\frac{\pi}{6}\right)\). 当 \(x\in[0,\pi]\) 时，\(y'>0\) 的区间为 \(0<x<\frac{\pi}{3}\) 或 \(\frac{5\pi}{6}<x<\pi\)，所以函数在 \([0,\frac{\pi}{3}]\) 和 \([\frac{5\pi}{6},\pi]\) 上单调递增.
\end{solution}

\begin{problem}
设甲，乙，丙三人每次射击命中目标的概率分别为 \(0.7\)，\(0.6\) 和 \(0.5\).

\begin{enumerate}
\item 三人各向目标射击一次，求至少有一人命中目标的概率及恰有两人命中目标的概率；
\item 若甲单独向目标射击三次，求他恰好命中两次的概率.
\end{enumerate}
\end{problem}

\begin{solution}
由于三人的射击相互独立，且甲、乙、丙每次命中的概率分别为 \(0.7\)、\(0.6\)、\(0.5\)，未命中的概率分别为 \(0.3\)、\(0.4\)、\(0.5\).

\begin{enumerate}
\item 三人都没有命中的概率为 \((1-0.7)(1-0.6)(1-0.5)=0.3\times0.4\times0.5=0.06\)，所以至少有一人命中的概率为 \(1-0.06=0.94\). 恰有两人命中的概率为
\(0.7\times0.6\times0.5+0.7\times0.4\times0.5+0.3\times0.6\times0.5=0.21+0.14+0.09=0.44\).
\item 甲每次命中的概率为 \(0.7\)，未命中的概率为 \(0.3\). 三次射击中恰好命中两次的概率为 \(\mathrm C_3^2(0.7)^2(0.3)=0.441\).
\end{enumerate}
\end{solution}

\begin{problem}
如图，四棱锥 \(P - ABCD\) 的底面是正方形，\(PA \perp\) 底面 \(ABCD\)，\(AE \perp PD\)，\(EF \parallel CD\)，\(AM = EF\).

\begin{enumerate}
\item 证明：\(MF\) 是异面直线 \(AB\) 与 \(PC\) 的公垂线；
\item 若 \(PA = 3AB\)，求直线 \(AC\) 与平面 \(EAM\) 所成角的正弦值.
\end{enumerate}

\FigureLayoutDeclare{img/2004/chongqing_liberal/q19_fig1.png}{6.0cm}{9bp 3bp 13bp 9bp}
\bitmapfigure[max width=0.36\linewidth]{img/2004/chongqing_liberal/q19_fig1.png}
\end{problem}

\begin{solution}
由图可知 \(E\in PD\)，\(F\in PC\)，\(M\in AB\). 因为底面 \(ABCD\) 是正方形，所以 \(AB\parallel CD\)，从而 \(AM\parallel EF\). 又已知 \(AM=EF\)，故四边形 \(AMFE\) 是平行四边形，于是 \(MF\parallel AE\).

由 \(PA\perp\) 平面 \(ABCD\)，且 \(AB\perp AD\)，得 \(AB\perp\) 平面 \(PAD\)，所以 \(AB\perp AE\)，进而 \(AB\perp MF\). 另外，\(DC\parallel AB\)，故 \(AE\perp DC\). 已知 \(AE\perp PD\)，而 \(PD\) 与 \(DC\) 是平面 \(PCD\) 内相交的两条直线，所以 \(AE\perp\) 平面 \(PCD\)，从而 \(AE\perp PC\). 由 \(MF\parallel AE\)，得 \(MF\perp PC\). 因此，\(MF\) 是异面直线 \(AB\) 与 \(PC\) 的公垂线.

设正方形边长 \(AB=a\)，则 \(PA=3a\). 建立直角坐标系，取 \(A\) 为原点，\(AB\)、\(AD\)、\(AP\) 分别为三个坐标轴正方向，则
\(A=(0,0,0)\)，\(B=(a,0,0)\)，\(D=(0,a,0)\)，\(C=(a,a,0)\)，\(P=(0,0,3a)\).

设 \(E=P+t(D-P)=(0,ta,3a(1-t))\). 由 \(AE\perp PD\)，得
\[
\overrightarrow{AE}\cdot\overrightarrow{PD}=(0,ta,3a(1-t))\cdot(0,a,-3a)=0
\]
即 \(t-9(1-t)=0\)，所以 \(t=\frac9{10}\). 因而 \(E=(0,\frac{9a}{10},\frac{3a}{10})\).

再设 \(F\in PC\). 由于 \(EF\parallel CD\)，可写 \(F=P+s(C-P)=(sa,sa,3a(1-s))\). 比较 \(E\)、\(F\) 的纵坐标，得 \(s=t=\frac9{10}\)，从而 \(EF=sa=ta=\frac{9a}{10}\). 由 \(AM=EF\) 且 \(M\in AB\)，得 \(M=(\frac{9a}{10},0,0)\).

平面 \(EAM\) 由向量 \(\overrightarrow{AM}=(\frac{9a}{10},0,0)\) 和 \(\overrightarrow{AE}=(0,\frac{9a}{10},\frac{3a}{10})\) 确定，因此它的一个法向量可取 \(\bs n=(0,-1,3)\). 又 \(\overrightarrow{AC}=(a,a,0)\)，设直线 \(AC\) 与平面 \(EAM\) 所成的角为 \(\theta\)，则
\[
\sin\theta=\frac{\left|\overrightarrow{AC}\cdot\bs n\right|}{\left|\overrightarrow{AC}\right|\left|\bs n\right|}=\frac{a}{a\sqrt{2}\cdot\sqrt{10}}=\frac{\sqrt{5}}{10}
\]
故所求正弦值为 \(\frac{\sqrt{5}}{10}\).
\end{solution}

\begin{problem}
某工厂生产某种产品，已知该产品的月生产量 \(x\)（\(\text{吨}\)）与每吨产品的价格 \(p\)（\(\text{元}/\text{吨}\)）之间的关系式为：\(p = 24200 - \frac{1}{5} x^2\)，且生产 \(x\text{ 吨}\) 的成本为 \(R = 50000 + 200x\)（\(\text{元}\)）. 问该产品每月生产多少吨产品才能使利润达到最大？最大利润是多少？（利润 \(=\) 收入-成本）
\end{problem}

\begin{solution}
设月利润为 \(\Pi(x)\)，则月收入为 \(xp=24200x-\frac15x^3\)，所以
\[
\Pi(x)=24200x-\frac15x^3-(50000+200x)=-\frac15x^3+24000x-50000
\]
其中 \(x\geqslant0\). 求导得 \(\Pi'(x)=24000-\frac35x^2\). 当 \(0\leqslant x<200\) 时，\(\Pi'(x)>0\)；当 \(x>200\) 时，\(\Pi'(x)<0\)，所以 \(x=200\) 时利润最大.

此时 \(\Pi(200)=-\frac15\times200^3+24000\times200-50000=3150000\). 因此，每月生产 \(200\text{ 吨}\) 时利润达到最大，最大利润为 \(3150000\text{ 元}\).
\end{solution}

\begin{problem}
设 \( p > 0 \) 是一常数，过点 \( Q(2p,0) \) 的直线与抛物线 \( y^{2} = 2px \) 交于相异两点 \(A\)，\(B\)，以线段 \(AB\) 为直径作圆 \(H\)（\(H\) 为圆心）. 试证抛物线顶点在圆 \(H\) 的圆周上，并求圆 \(H\) 的面积最小时直线 \(AB\) 的方程.
\end{problem}

\begin{solution}
记抛物线顶点为 \(O\). 因为 \(x\) 轴 \(y=0\) 与抛物线只交于顶点 \(O\)，不能与抛物线交于相异两点，所以所取直线不是 \(x\) 轴，可设其方程为 \(x=2p+\lambda y\). 设它与抛物线的交点为 \(A=(x_1,y_1)\)、\(B=(x_2,y_2)\)，则 \(y_1\)，\(y_2\) 是方程
\[
y^2-2p\lambda y-4p^2=0
\]
的两个不相等的实根. 因而 \(y_1+y_2=2p\lambda\)，\(y_1y_2=-4p^2\)，且 \(x_i=2p+\lambda y_i\).

于是
\[
x_1x_2=(2p+\lambda y_1)(2p+\lambda y_2)=4p^2+2p\lambda(y_1+y_2)+\lambda^2y_1y_2=4p^2
\]
所以 \(\overrightarrow{OA}\cdot\overrightarrow{OB}=x_1x_2+y_1y_2=0\)，由勾股定理的逆定理，\(OA\perp OB\). 因此顶点 \(O\) 在以 \(AB\) 为直径的圆 \(H\) 上.

又由根的差的平方等于判别式，得 \((y_1-y_2)^2=4p^2(\lambda^2+4)\). 因为 \(x_1-x_2=\lambda(y_1-y_2)\)，所以
\[
AB^2=(x_1-x_2)^2+(y_1-y_2)^2=4p^2(\lambda^2+1)(\lambda^2+4)
\]
圆 \(H\) 的半径为 \(AB/2\)，故其面积为
\[
S=\pi\left(\frac{AB}{2}\right)^2=\pi p^2(\lambda^2+1)(\lambda^2+4)
\]
令 \(u=\lambda^2\geqslant0\)，则 \((u+1)(u+4)=u^2+5u+4\geqslant4\)，等号当且仅当 \(u=0\)，即 \(\lambda=0\). 此时直线 \(AB\) 的方程为 \(x=2p\)，圆 \(H\) 的面积最小.
\end{solution}

\begin{problem}
设数列 \(\{a_n\}\) 满足：\(a_1 = 2\)，\(a_2 = \frac{5}{3}\)，\(a_{n+2} = \frac{5}{3} a_{n+1} + \frac{1}{3} a_n\)（\(n = 1\)，\(2\)，\(\cdots\)）.

\begin{enumerate}
\item 令 \(b_n = a_{n+1} - a_n\)（\(n = 1\)，\(2\)，\(\cdots\)），求数列 \(\{b_n\}\) 的通项公式；
\item 求数列 \(\{na_n\}\) 的前 \(n\) 项和 \(S_n\).
\end{enumerate}
\end{problem}

\begin{solution}
设特征方程为 \(r^2-\frac53r-\frac13=0\)，其两个根为
\(\alpha=\frac{5+\sqrt{37}}{6}\)，\(\beta=\frac{5-\sqrt{37}}{6}\).
由 \(\alpha+\beta=\frac53\)、\(\alpha\beta=-\frac13\)，以及初值 \(a_1=2=1+1\)、\(a_2=\frac53=\alpha+\beta\)，可得
\[
a_n=\alpha^{n-1}+\beta^{n-1}
\]

\begin{enumerate}
\item 由 \(b_n=a_{n+1}-a_n\)，得
\[
b_n=(\alpha-1)\alpha^{n-1}+(\beta-1)\beta^{n-1}=\frac{\sqrt{37}-1}{6}\left(\frac{5+\sqrt{37}}{6}\right)^{n-1}-\frac{\sqrt{37}+1}{6}\left(\frac{5-\sqrt{37}}{6}\right)^{n-1}
\]
\item 由上式，
\[
S_n=\sum_{k=1}^{n}ka_k=\sum_{k=1}^{n}k\left(\alpha^{k-1}+\beta^{k-1}\right)
\]
对任意 \(r\ne1\)，由等比数列求和公式求导可得
\[
\sum_{k=1}^{n}kr^{k-1}=\frac{\mathrm d}{\mathrm dr}\left(\frac{1-r^{n+1}}{1-r}\right)=\frac{1-(n+1)r^n+nr^{n+1}}{(1-r)^2}
\]
因此
\[
S_n=\frac{1-(n+1)\alpha^n+n\alpha^{n+1}}{(1-\alpha)^2}+\frac{1-(n+1)\beta^n+n\beta^{n+1}}{(1-\beta)^2}
\]
其中 \(\alpha=\frac{5+\sqrt{37}}{6}\)，\(\beta=\frac{5-\sqrt{37}}{6}\).
\end{enumerate}
\end{solution}
